% ==================================================================
% id:       seq-005
% topic:    סדרות
% subtopic: הגדרת הגבול (ε, N)
% solution: full
% ==================================================================
\textbf{תהי
$(a_n)_{n=1}^{\infty}$
סדרה
ויהי
$L\in\rr$.}
\bE
\item
    הראו כי אם
    $a_n\rightarrow L$
    אז
    $|a_n|\rightarrow |L|$.
\hint{השתמשו באי-שוויון המשולש ההפוך $||x|-|y||\le |x-y|$.}
    \begin{solbox}
     יהי
    $\varepsilon>0$.
    מהגדרת הגבול של
    $a_n\to L$,
    נובע כי
    קיים
    $N$
    כך ש-
    $|a_n-L|<\varepsilon$
    לכל
    $n>N$.
    מאי-שיוויון המשולש ההפוך נסיק כי
    $$||a_n|-|L||\leq|a_n-L|<\varepsilon$$
    לכל
    $n>N$.
    לסיכום, הראנו כי לכל
    $\varepsilon>0$
    קיים
    $N$
    טבעי כך שלכל
    $n>N$
    מתקיים
    $$\left| |a_n|-|L|\right|<\varepsilon$$
    שזוהי בדיוק ההגדרה של
    $|a_n|\to |L|$.

\end{solbox}
    \item
    הראו כי הכיוון ההפוך אינו נכון:
    מצאו
    $a_n$
    המקיימת
    $a_n\not\rightarrow L$
    אבל
    $|a_n|\rightarrow |L|$.
    
\hint{חשבו על סדרה שהערך המוחלט שלה מתכנס אבל הסימן שלה משתנה.}
\begin{solbox}
    הסדרה
    $a_n=(-1)^n$
    מקיימת
    $|a_n|=1\to 1=L$
    אבל כמובן אינה מתכנסת ולכן
    $a_n\not\to 1$.
    שימו לב כי גם
    $L=-1$
    היה עובד כאן.
\end{solbox}
    \item
    הראו כי אם
    $a_n\rightarrow L$
    אז
    $a_n^2\rightarrow L^2$.
    \begin{solbox}
    זהו שימוש ישיר באריתמטיקה של כפל סדרות מתכנסות: מתקיים
    $a_n\to L$
    וגם
    $a_n\to L$,
    לכן
    $a_n^2=a_n\cdot a_n\to L\cdot L=L^2$.
\end{solbox}
      \item
    הראו כי הכיוון ההפוך אינו נכון: מצאו
    $a_n$
    המקיימת
    $a_n\not\rightarrow L$
    אבל
    $a_n^2\rightarrow L^2$.
    
\hint{נסו לחשוב על סדרה שלא מתכנסת, אבל העלאה שלה בריבוע „מסדרת” אותה והופכת אותה למתכנסת.}
\begin{solbox}
    גם כאן הדוגמה של סעיף (ב) היא רלוונטית וטובה:
    עבור
    $a_n=(-1)^n$,
    מתקיים
    $a_n^2=1\to 1=1^2$
    אבל
    $a_n\not\to 1$.
\end{solbox}
    \eE
